% Part: lambda-calculus
% Chapter: syntax
% Section: abbreviated-syntax

\documentclass[../../../include/open-logic-section]{subfiles}

\begin{document}

\olfileid{lam}{syn}{abb}
\olsection{د نحو لنډې نښې}

د \olref[trm]{defn:term} له مخې تعريف شوي ترمونه کله کله په ليکلو
کښې درانۀ وي؛ نو د نحو د لا لنډې بڼې ټاکل ګټور دي۔ البته بايد پام
وکړو چې په لنډه نښه کښې ترمونه هم په يکتا توګه ولوستل شي۔ يوه ډېره
کارېدونکې بڼه، چې \emph{لنډ شوي ترمونه} بلل کېږي، داسې ده:

\begin{enumerate}
\item چې قوسونه پرېښودل شي، تطبيق له کيڼ نه ښي ته کېږي۔ مثلاً، که
  $M$، $N$، $P$ او~$Q$ ترمونه وي، نو $MNPQ$ د $(((MN)P)Q)$
  لنډه نښه ده۔
\item همداسې، چې قوسونه پرېښودل شي، د لامبډا تجريد د امکان تر
  ټولو پراخه ساحه اخلي۔ مثلاً، $\lambd[x][MNP]$ د
  $(\lambd[x][MNP])$ په توګه لوستل کېږي۔
\item يوه لامبډا د څو متغيرونو د تجريد دپاره کارېدے شي۔ مثلاً،
  $\lambd[xyz][M]$ د $\lambd[x][\lambd[y][\lambd[z][M]]]$
  لنډه نښه ده۔
\end{enumerate}

د بېلګې په توګه،
\[
\lambd[xy][xxyx \lambd[z][xz]]
\]
د دې لنډه نښه ده:
\[
(\lambd[x][(\lambd[y][((((xx)y)x)(\lambd[z][(xz)]))])]).
\]

\begin{prob}
لنډ شوے ترم $\lambd[g][(\lambd[x][g (x x)])
  \lambd[x][g (x x)]]$ په بشپړه بڼه وليکئ۔
\end{prob}

\end{document}
